%!TeX root=MemoriaTFG.tex

\chapter{Estado del arte}
\label{cap:sota}

\section{Aprendizaje profundo aplicado a dermatología}
\label{sec:dl-dermatologia}

El trabajo de Esteva et~al.~\cite{esteva2017} establece la primera
referencia consolidada de clasificación dermatológica con redes
convolucionales, con rendimiento equivalente al de un panel de
dermatólogos certificados sobre clasificación binaria
maligno/benigno. Las extensiones posteriores amplían el espectro
diagnóstico y la diversidad de datos~\cite{tschandl2018,ham10000} y
formulan tareas adicionales como segmentación binaria de
lesión~\cite{isic2018} y clasificación binaria
melanoma/no-melanoma con anotaciones estructuradas
multi-criterio~\cite{kawahara2019}. El paradigma común permanece
estable: entrenamiento desde pesos de ImageNet, una arquitectura por
tarea y dependencia de anotación experta. Tres limitaciones lo
acotan: el coste de la anotación clínica, la fragmentación
taxonómica entre conjuntos públicos y la pérdida de rendimiento
fuera del dominio de entrenamiento, particularmente acusada en
fototipos cutáneos elevados~\cite{fitzpatrick17k,ddi}.

\section{Inteligencia artificial en imagen dermatoscópica y clínica}
\label{sec:sota-modalidades}

La modalidad en la que la inteligencia artificial está más madura es
la dermatoscopia aplicada al melanoma. El motivo es, en buena medida,
la disponibilidad de datasets públicos amplios y bien anotados y de
una serie de competiciones internacionales que han fijado un banco de
pruebas común. Los retos de la \emph{International Skin Imaging
Collaboration} (ISIC), celebrados de forma anual desde 2016 junto a
congresos de visión por computador, se apoyan en un archivo público
de más de medio millón de imágenes de lesiones cutáneas, en su
mayoría dermatoscópicas~\cite{isicchallenge}. Sobre datasets de este tipo,
como HAM10000~\cite{ham10000}, los modelos alcanzan un rendimiento
comparable al del dermatólogo experto en la distinción entre melanoma
y lesiones benignas; una revisión sistemática con meta-análisis
reciente concluye que, para la clasificación de cáncer de piel, los
algoritmos igualan a los dermatólogos expertos y superan a los
clínicos no especializados~\cite{salinas2024}.

Fuera del melanoma y del cáncer cutáneo queratinocítico, el avance es
bastante menor. La mayoría de los datasets dermatoscópicos y de la
literatura de evaluación se concentran en la detección de cáncer de
piel ---una revisión sistemática de $232$ estudios circunscribe el
grueso de las aplicaciones de IA en dermatología al melanoma y al
carcinoma queratinocítico~\cite{jairath2024}---, y la extensión a
tareas multiclase sobre el conjunto amplio de
dermatosis es más reciente y está mucho menos consolidada. La propia
tarea de segmentación de la lesión se ha formulado también sobre
datasets dermatoscópicos de cáncer~\cite{isic2018}, lo que ilustra
hasta qué punto la modalidad gravita en torno al problema oncológico.

La fotografía clínica o macroscópica ---la obtenida con una cámara
convencional o un teléfono, sin lente de contacto ni aumento--- está
mucho menos explorada. El trabajo de Esteva et~al.~\cite{esteva2017}
mostró que una red entrenada sobre imágenes clínicas podía clasificar
cáncer de piel a nivel de dermatólogo, y sistemas posteriores como el
de Liu et~al.~\cite{liu2020} ampliaron el alcance a las $26$
dermatosis más frecuentes en atención primaria a partir de casos de
teledermatología. Aun así, estos esfuerzos son escasos frente al
volumen del trabajo dermatoscópico, se concentran en un número
reducido de sistemas y siguen condicionados por problemas de
generalización: el rendimiento cae sobre los fototipos cutáneos más
oscuros y sobre los diagnósticos menos
frecuentes~\cite{ddi,fitzpatrick17k,tjiu2025}.

Esta asimetría delimita el terreno menos trabajado. La dermatoscopia
para melanoma está relativamente resuelta, pero el resto del espectro
de patologías ---y, en especial, la fotografía clínica--- sigue siendo
poco explorado, y casi todo el material disponible está en inglés.
Incluso los recursos más recientes que enriquecen imágenes con
descripciones clínicas, como SkinCAP~\cite{skincap2024}, se construyen
sobre datasets clínicos en inglés ya existentes. La parte menos
explorada de este panorama es, por tanto, la fotografía clínica en
lenguas distintas del inglés y con texto narrativo de caso real.

\section{Modelos fundacionales: arquitectura y objetivos}
\label{sec:fundacionales-base}

Los modelos fundacionales sustituyen el entrenamiento específico por
encoders preentrenados de forma auto-supervisada o multimodal sobre
datasets visuales masivos. La arquitectura predominante es el Vision
Transformer~\cite{vit2021}, que procesa la imagen como una secuencia
de parches y aprende dependencias globales mediante mecanismos de
atención. Los objetivos de preentrenamiento más extendidos son tres:
el \emph{masked image modeling}, en el que la red reconstruye
parches ocultos a partir del contexto visible; la destilación
auto-supervisada entre vistas; y la alineación contrastiva
imagen-texto al estilo CLIP~\cite{clip2021}. La adaptación posterior
a tareas específicas opera por linear probe sobre representaciones
congeladas, fine-tuning supervisado, o transferencia directa
zero-shot guiada por \emph{prompts} en los modelos con rama
lingüística.

\section{Encoders visuales auto-supervisados}
\label{sec:encoders-visuales}

DINOv2~\cite{dinov2} es la referencia abierta de encoder visual
auto-supervisado de propósito general, entrenado sobre
$\sim 142$ millones de imágenes naturales sin etiquetas mediante
destilación entre vistas. Sus representaciones transfieren a tareas
de clasificación, recuperación y segmentación sobre dominios
naturales y biomédicos con ajuste mínimo, y constituyen una base
habitual para adaptación a dominios especializados.

En el dominio dermatológico, PanDerm~\cite{panderm2025} es un
encoder visual ViT (Base $\sim 86$\,M y Large $\sim 307$\,M
parámetros) entrenado sobre $\sim 2{,}1$ millones de imágenes
dermatológicas en cuatro modalidades (fotografía corporal total,
dermatoscopia, fotografía clínica e histopatología). El
preentrenamiento utiliza \emph{masked latent modeling}: la red
reconstruye representaciones latentes de regiones ocultas en lugar
de píxeles, alineándolas con las que produce un codificador CLIP-Large
congelado como objetivo (\emph{teacher}) ---del que solo se usa la
rama de imagen, de modo que PanDerm sigue siendo un modelo de
imagen---. Sus autores reportan una mejora media del
$+10{,}2\%$ sobre el estado del arte previo en $28$ tareas
dermatológicas. La disponibilidad simultánea de pesos, código y
benchmarks lo convierte en el primer modelo fundacional dermatológico
verificable de forma independiente.

\section{Modelos vision-language contrastivos}
\label{sec:vision-language}

CLIP~\cite{clip2021} establece el preentrenamiento contrastivo
imagen-texto sobre $400$ millones de pares de la web. El objetivo
InfoNCE aproxima la imagen a su descripción dentro del lote y la
aleja del resto. El espacio compartido habilita dos modos
operativos: clasificación zero-shot guiada por \emph{prompts}
textuales (la clase ganadora es la de mayor similitud coseno entre
el vector visual y los vectores textuales de las clases candidatas)
y recuperación bidireccional imagen-texto sobre datasets
precomputados.

SigLIP~\cite{siglip2023} sustituye la pérdida softmax de CLIP por
una pérdida sigmoide que reduce la dependencia del tamaño de batch y
escala con mejor estabilidad. Su variante Large
(SO400M) dispone de embedding de $1\,152$ dimensiones y
$\sim 878$\,M parámetros. BiomedCLIP~\cite{biomedclip} extiende
el paradigma al dominio biomédico sobre PMC-15M
($\sim 15$ millones de pares figura-leyenda de literatura
científica); su cobertura dermatológica específica es limitada por
el sesgo de selección del dataset.

En dermatología, Derm1M y la familia DermLIP~\cite{derm1m} añaden la
rama lingüística contrastiva al estado del arte de la disciplina.
Derm1M reúne $1\,029\,761$ pares imagen-texto sobre
$403\,563$ imágenes únicas, organizados según una ontología
jerárquica de cuatro niveles construida por dermatólogos. La familia
DermLIP combina varios encoders visuales (PanDerm Base,
ConvNeXt-Large) con varios encoders textuales (PubMedBERT extendido
a $256$ tokens, GPT-2 dermatológico con tokenizer de CLIP a
$77$ tokens). La configuración con mejor rendimiento reportada
por los autores combina PanDerm Base como encoder visual y
PubMedBERT como encoder textual, superando a BiomedCLIP en los
benchmarks zero-shot evaluados.

La implementación operativa de la recuperación visual densa sobre
datasets a gran escala requiere índices vectoriales optimizados para
búsqueda aproximada del vecino más próximo. FAISS~\cite{faiss2019}
y estructuras basadas en grafos del tipo Hierarchical Navigable
Small Worlds permiten consultas en el orden de $10^2$
milisegundos sobre datasets de centenares de miles a millones de
\emph{embeddings} precomputados.

\section{Modelos integrados con zero-shot nativo}
\label{sec:zero-shot-nativo}

DermFM-Zero~\cite{dermfmzero2026} integra el preentrenamiento visual
auto-supervisado y la alineación contrastiva imagen-texto en un
único modelo con capacidad zero-shot nativa. El entrenamiento se
estructura en dos fases: \emph{masked latent modeling} sobre tres
millones de imágenes en la primera, y \emph{bootstrapped contrastive
learning} sobre un millón de pares texto-imagen en la segunda. El
encoder textual es PubMedBERT. El modelo incorpora interpretabilidad
emergente mediante Sparse Autoencoders~\cite{templeton2024}, redes
auxiliares que descomponen la representación interna en un
diccionario disperso de direcciones asociables a conceptos
clínicamente identificables. Sus autores reportan validación en
tres \emph{reader studies} con más de mil cien clínicos. Sus pesos
no son públicos en la fecha de cierre del presente trabajo, lo que
restringe su tratamiento a la discusión conceptual.

BLIP-2~\cite{blip2_2023} introduce un patrón arquitectónico
distinto: conecta un encoder visual congelado con un modelo de
lenguaje preentrenado mediante un módulo intermedio ligero
(\emph{Q-Former} o \emph{Querying Transformer}) que transforma las
representaciones visuales en una secuencia reducida de tokens
interpretables por el modelo de lenguaje. Sólo el Q-Former se
entrena. No existe, en la fecha de cierre, una adaptación pública
específica de BLIP-2 al dominio dermatológico.

\section{Segmentación promptable}
\label{sec:segmentacion-sota}

SAM~\cite{sam2023} establece el paradigma de segmentación
promptable: produce máscaras a partir de \emph{prompts} visuales
(puntos, cajas, máscaras parciales) sobre imagen natural. Su
preentrenamiento utiliza $\sim 11$ millones de imágenes con
más de mil millones de máscaras generadas mediante un proceso
semi-automatizado. SAM2~\cite{sam2} amplía el enfoque al dominio
del vídeo y mejora el rendimiento sobre imágenes estáticas.

La transferencia a dominios médicos requiere adaptación supervisada
sobre conjuntos con máscaras binarias. La estrategia habitual
combina un modelo SAM/SAM2 preentrenado con fine-tuning sobre
ISIC2018~\cite{isic2018} y adaptadores de bajo rango
(\emph{Low-Rank Adaptation}, LoRA)~\cite{lora}, que
incorporan un número reducido de parámetros entrenables manteniendo
congelado el peso original del encoder visual. Esta estrategia
reduce los requisitos de memoria y cómputo y permite la adaptación
sobre hardware modesto.

\section{Modelos de lenguaje multimodales generalistas}
\label{sec:llm-multimodal}

Los modelos de lenguaje multimodales preentrenados sobre datasets web
masivos constituyen una referencia de comparación pertinente: GPT-4o
(OpenAI), Gemini~2.5~Pro (Google DeepMind) y la familia MedGemma
(versiones 4B Vision y 27B Text, Google). Su número de parámetros
supera al de los encoders fundacionales específicos en uno o dos
órdenes de magnitud. Su modo de uso característico es la inferencia
guiada por \emph{prompts} sin fine-tuning específico por
tarea. Trabajos recientes documentan su uso conversacional para
soporte al razonamiento clínico~\cite{amie2024,deepmind_aicoclinician}.
Los modelos de acceso restringido (GPT-4o, Gemini) sólo se evalúan
mediante API; MedGemma libera pesos y se evalúa en local.

\section{Problemas metodológicos y guías de reporte}
\label{sec:problemas-metodologicos}

La aplicación de modelos fundacionales a imagen médica plantea tres
problemas metodológicos directamente relevantes para el presente
trabajo. El primero es la fuga de información entre conjuntos de
evaluación y dataset de preentrenamiento. La construcción de Derm1M
sobre fuentes web abiertas implica que datasets como Dermnet, HIBA
o MSKCC pueden solaparse con su dataset de entrenamiento, lo que
sesga al alza las métricas zero-shot reportadas. El segundo es la
heterogeneidad taxonómica entre datasets, que impide comparaciones
directas sin un paso explícito de armonización. El tercero es la
disparidad de rendimiento por fototipo cutáneo, documentada en
Fitzpatrick17k~\cite{fitzpatrick17k} y atendida específicamente por
SkinCon~\cite{skincon} mediante anotación basada en $48$
conceptos clínicos visuales.

Las guías de reporte recientes para inteligencia artificial en
imagen médica ---CLAIM~\cite{claim2024} y la actualización
TRIPOD+AI~\cite{tripodai2024}--- establecen la transparencia metodológica como
condición de defensibilidad: cohorte, particiones, métricas,
incertidumbre y limitaciones deben constar explícitamente en el
informe completo. El reporte de pruebas estadísticas (intervalos de
confianza por \emph{bootstrap}, comparaciones múltiples,
contrastes pareados de AUROC mediante el test de
DeLong~\cite{delong1988}, y tamaños muestrales por subgrupo)
constituye un requisito creciente en publicaciones de la disciplina.

\section{Resumen del estado del arte}
\label{sec:resumen-sota}

El recorrido anterior puede resumirse en una idea. El aprendizaje
profundo alcanzó un rendimiento comparable al del especialista en la
detección de melanoma~\cite{esteva2017}, pero al precio de una
arquitectura por tarea, una fuerte dependencia de anotación experta y
una caída notable de rendimiento fuera del dominio de entrenamiento,
especialmente sobre los fototipos cutáneos más
oscuros~\cite{fitzpatrick17k,ddi}. Los \emph{modelos fundacionales}
cambiaron esa lógica: un mismo encoder, preentrenado a gran
escala, se reutiliza en muchas tareas mediante adaptación ligera. En
dermatología, las referencias abiertas son PanDerm ---un encoder
visual entrenado sobre las cuatro modalidades de la
especialidad~\cite{panderm2025}--- y la familia DermLIP, que añade una
rama de texto y la clasificación \emph{zero-shot}~\cite{derm1m}; el
modelo más reciente con \emph{zero-shot} nativo, DermFM-Zero, no
libera sus pesos~\cite{dermfmzero2026}. Alrededor de ellos conviven
encoders generalistas (DINOv2), modelos contrastivos
imagen-texto (CLIP, SigLIP, BiomedCLIP), la segmentación
\emph{promptable} de la familia SAM~\cite{sam2} y los grandes modelos
de lenguaje multimodales como comparadores de contexto.

Tres problemas atraviesan toda esta literatura y condicionan
cualquier lectura de sus resultados: la fuga de información entre los
datasets de evaluación y los de preentrenamiento, que infla las
cifras; la incompatibilidad entre las taxonomías de cada dataset, que
impide compararlos directamente; y la desigualdad de rendimiento
entre fototipos cutáneos. A ello se suma una exigencia creciente de
transparencia metodológica ---cohorte, particiones, incertidumbre y
contrastes estadísticos--- recogida en guías como
TRIPOD+AI~\cite{tripodai2024}.

De todo lo anterior se desprenden las carencias que motivan el trabajo.
Cada modelo se publica evaluado por sus propios autores, sobre sus
propios \emph{benchmarks} y con protocolos distintos: falta una
comparación independiente, bajo un mismo protocolo y sobre datasets
externos comunes, realizada por un equipo ajeno a quien creó los
modelos. Casi todo el material y los encoders de texto están en
inglés y centrados en dermatoscopia; la fotografía clínica en
español, con texto narrativo de caso real, es prácticamente
inexistente ---los recursos recientes de \emph{captions} clínicas
ricas, como SkinCAP~\cite{skincap2024}, enriquecen datasets en inglés
ya existentes, pero no cubren esa carencia---. Y entre la evaluación
sobre \emph{benchmarks} y el despliegue real en un hospital media una
distancia operativa que los trabajos publicados no abordan.

De este panorama se desprende el plan de este trabajo. El primer paso
es reproducir y validar PanDerm ---el modelo fundacional dermatológico
abierto más completo--- sobre un conjunto común de datasets externos,
comprobando que su comportamiento se corresponde con el publicado y
comparándolo de forma homogénea con los demás modelos disponibles;
ello exige armonizar las taxonomías heterogéneas de los datasets y
auditar su posible solapamiento con los datos de preentrenamiento.
Sobre esta base se quiere explorar un conjunto de ideas
que el propio estado del arte ---y en particular DermFM-Zero--- apunta
como prometedoras: el trabajo en el espacio latente y el alineamiento
contrastivo imagen-texto al estilo CLIP, la interpretabilidad mediante
\emph{sparse autoencoders} ---adaptada al español para explicar las
imágenes--- y la clasificación \emph{zero-shot} guiada por lenguaje.
El horizonte de todo ello es el reconocimiento sobre fotografía
clínica ---en clasificación y en segmentación---, la modalidad menos
trabajada y de mayor interés asistencial.
